% !TEX TS-program = pdflatex
% !TEX encoding = UTF-8 Unicode
% !BIB TS-program = bibtex
%
% ============================================================================
% paper.tex -- Standard compliance bounds origin--destination identifiability
%              in GBFS bike-sharing feeds: a cost-learning and sampling-horizon
%              analysis.
%
% Target  : Computer Standards & Interfaces (Elsevier).
% Authors : R. Fosse and G. Pallares, CESI LINEACT (EA 7527), Montpellier, 2026.
%
% Structure: linear engineering/OR flow (Elsevier/CS&I), NOT Nature/Science --
%   Intro -> Related Work (front) -> Formulation -> Theory (identifiability,
%   bias) -> Sampling-horizon law -> Empirical audit (method folded in) ->
%   Discussion/policy -> Conclusion. Long proofs live in the SI.
%
% >>> SUBMISSION NOTE <<<
%   Scaffolded from the org's IOP `paper-template`; ships `iopjournal.cls`.
%   CS&I is an ELSEVIER journal: before submission switch to
%       \documentclass[review,3p,times]{elsarticle}
%   and migrate the IOP back-matter macros (\ack, \funding, \roles, \data)
%   into an elsarticle \begin{frontmatter}...\end{frontmatter} + CRediT block.
%   The body, theorems and bibliography are class-agnostic. See README.md.
%
% Compile :
%   pdflatex paper.tex && bibtex paper && pdflatex paper.tex && pdflatex paper.tex
% ============================================================================
\documentclass{iopjournal}      % add [anonymous] for double-anonymous review

\usepackage[T1]{fontenc}
\usepackage[utf8]{inputenc}
\usepackage{microtype}
\usepackage{csquotes}
\usepackage{amssymb,amsmath,amstext,amsthm}

\usepackage{graphicx}
\graphicspath{{figures/}}

\usepackage{booktabs}
\usepackage{array}
\usepackage{float}
\usepackage{enumitem}
\setlist{topsep=2pt,itemsep=2pt,parsep=0pt}
\usepackage{url}

\usepackage[authoryear,round]{natbib}
\bibliographystyle{plainnat}

\IfFileExists{orcid.pdf}{}{%
  \renewcommand{\orcid}[1]{\textsuperscript{\href{https://orcid.org/##1}{\textsf{[ORCID]}}}}%
}

% Target journal (rendered in the banner).  See SUBMISSION NOTE above.
\newcommand{\journalname}{Computer Standards \& Interfaces}
\newcommand{\journalshort}{Comput.\ Stand.\ Interfaces}

\usepackage{fancyhdr}
\fancyhead[L]{{\small \sf Preprint --- working draft}\hspace{5mm}%
              {\it \journalshort}\ (submitted)}
\fancyhead[R]{Foss\'e {\it et al}\ }
\fancyhead[C]{}

\widowpenalty=10000
\clubpenalty=10000
\displaywidowpenalty=10000
\raggedbottom

\makeatletter
\renewcommand{\articletype}[1]{%
   {\vspace*{-8mm}\noindent \Large \sf \journalname}%
   \vspace*{8mm} \noindent\reversemarginpar
   \marginpar{\vspace{-3mm} {\color{gray}\hrule} \ \\ Preprint\\
       {\color{gray}\hrule} \ \\
       \tiny {\sf VERSION} {\small \\ v0.2 draft}\\ \\
       {\sf TARGET} {\small \\ \journalshort}}%
   {\scriptsize \sf{\bfseries \MakeUppercase{#1}}}}
\makeatother

% Theorem environments (shared counter: Thm 1, Def 2, Prop 3, Cor 4, Thm 5, Cor 6)
\theoremstyle{plain}
\newtheorem{theorem}{Theorem}
\newtheorem{proposition}[theorem]{Proposition}
\newtheorem{lemma}[theorem]{Lemma}
\newtheorem{corollary}[theorem]{Corollary}
\theoremstyle{definition}
\newtheorem{definition}[theorem]{Definition}
\theoremstyle{remark}
\newtheorem{remark}[theorem]{Remark}
\newtheorem{conjecture}[theorem]{Conjecture}

% Math shortcuts
\newcommand{\E}{\mathbb{E}}
\newcommand{\R}{\mathbb{R}}
\newcommand{\Ncal}{\mathcal{N}}              % separable null space {f_i + g_j}
\newcommand{\Pperp}{\Pi_{\Ncal^\perp}}        % projection onto the identifiable part
\newcommand{\KL}{\mathrm{KL}}
\newcommand{\OD}{\mathrm{OD}}
\newcommand{\neff}{n_{\mathrm{eff}}}
\newcommand{\Tstar}{T^{\star}}
\newcommand{\Pstar}{P^{\star}}
\newcommand{\cstar}{c^{\star}}

% Drafting marker -- remove all uses before v1.0.
\newcommand{\Todo}[1]{\textbf{[To confirm: #1]}}

% Figure-or-placeholder: includes the PDF if rendered, else a compile-safe box.
% Usage: \figorplaceholder{basename}{width}{placeholder caption text}
\newcommand{\figorplaceholder}[3]{%
  \IfFileExists{figures/#1.pdf}{\includegraphics[width=#2]{#1}}{%
    \fbox{\parbox[c][3.2cm][c]{#2}{\centering\footnotesize\Todo{#3}}}}}

% =========================================================================
\begin{document}

\articletype{Paper}

\title{Standard compliance bounds origin--destination identifiability in
GBFS bike-sharing feeds: a cost-learning and sampling-horizon analysis}

\author{Rohan Foss\'e$^{1,*}$\orcid{0009-0002-2195-0198},
        Ga\"el Pallares$^{2}$\orcid{0009-0002-8680-604X}}

\affil{$^{1}$ CESI \'Ecole d'Ing\'enieurs, Montpellier, France
       \quad$^{2}$ CESI LINEACT (EA 7527), Montpellier, France\\[0.2em]
       $^{*}$ Author to whom any correspondence should be addressed.}

\email{rfosse@cesi.fr}

\keywords{GBFS; origin--destination reconstruction; data-standard compliance;
entropic optimal transport; shared micromobility}

\begin{abstract}
Reconstructing origin--destination (OD) flows from public bike-sharing feeds
is split between two regimes: macroscopic inference from station stock
counts, which is underdetermined, and deterministic tracking of per-vehicle
identifiers, which is exposed to temporal aliasing and to the
privacy-driven rotation of identifiers recommended by the GBFS standard.
We recast the problem as a hybrid in which a partial, identifier-tracked
micro-sample supplies the deterrence (cost) structure of an entropic
optimal-transport / gravity model whose margins are pinned by the station
counts. Two results follow. First, the transport cost is identifiable only
modulo additive station effects; the bias induced by incomplete observation
equals the \emph{non-separable} component of the log-selection probability,
so station-emptiness censoring cancels exactly, whereas 60-second
polling---whose capture probability is a function of travel time, the very
quantity being estimated---does not. Second, the estimator yields a
\emph{sampling-horizon law}: the continuous-collection time needed to
reconstruct OD to target accuracy $\delta$ scales as $\delta^{-4}$ and
inversely with the identifier-persistence rate $q$, implying a
standard-compliance threshold below which operational OD monitoring is
infeasible. An audit of \Todo{$N$} French free-floating networks instantiates
the bound and exhibits a sharp bimodal split between standard-compliant
(identifier-rotating, untrackable) and non-compliant (persistent, trackable)
feeds. The contribution is therefore a \emph{normative feasibility audit}---a
characterisation of when standard compliance makes OD recovery impossible in
operational time---and not a new reconstruction estimator: the estimator's
convergence is established on synthetic ground truth, while the real feeds serve
only to instantiate the bound's constants and place each network on the
feasibility map.
\end{abstract}

% =========================================================================
\section{Introduction}
\label{sec:intro}

Origin--destination (OD) matrices are the primitive of every operational and
planning question in shared micromobility: fleet rebalancing, station siting,
demand forecasting, and the equity audits that increasingly condition public
funding all consume an OD matrix as input. Yet the data that operators expose
publicly was never designed to carry it. The General Bikeshare Feed
Specification (GBFS), the de-facto open standard adopted by well over a
thousand systems worldwide, is a \emph{real-time availability} interface: it
advertises which vehicles are where, right now, and explicitly does not record
what happened between two states. OD must therefore be \emph{reconstructed},
and the literature is cleaved along the only two routes the data permits.

The first route is macroscopic: treat the per-station stock changes as link or
node counts and invert. This is the classical OD-from-counts problem, and it is
mathematically mature---but it is also \emph{underdetermined}: many matrices
reproduce the same observed marginals, so the interior of the OD matrix is
fixed by a prior, not by the data. The second route is microscopic: when a feed
exposes a persistent per-vehicle identifier, individual vehicles can be tracked
between snapshots and their trips read off directly. This route is exact in
principle, but fragile in practice for two reasons that the literature treats as
nuisances rather than first-class objects. It is degraded by \emph{temporal
aliasing}, because polling at a finite cadence misses short or fast trips; and,
more fundamentally, it is \emph{deliberately defeated by the standard itself},
whose privacy guidance recommends rotating the vehicle identifier after every
trip precisely to prevent the trip reconstruction the method relies on.

This places the analyst in a paradox that, to our knowledge, has never been made
quantitative: the very property that makes a feed useful for OD reconstruction
(identifier persistence) is the property the standard asks operators to remove.
We argue that the feasibility of OD reconstruction is therefore not a modelling
choice but a \emph{function of the feed's compliance with the standard}, and we
make that dependence precise.

\paragraph{Contributions.}
We (i) recast OD reconstruction as a hybrid in which an identifier-tracked
micro-sample supplies the cost structure of an entropic optimal-transport
(equivalently, maximum-entropy gravity) model whose margins are pinned by the
station counts (Section~\ref{sec:formulation}); (ii) prove that the learned cost
is identifiable only modulo additive station effects, and decompose the
observation bias into the \emph{non-separable} part of the log-selection
probability, which shows that station-emptiness censoring cancels while
60-second aliasing does not (Section~\ref{sec:theory}); (iii) derive a
sampling-horizon law $\Tstar \propto \delta^{-4} q^{-1}$ and the
standard-compliance threshold it implies, and identify cross-network pooling as
the only lever that keeps the horizon finite (Section~\ref{sec:horizon}); and
(iv) instantiate the bound through a feed-by-feed audit of \Todo{$N$} French
networks, whose identifier-persistence rates split sharply into compliant and
non-compliant regimes (Section~\ref{sec:audit}). Section~\ref{sec:discussion}
translates the threshold into a concrete recommendation for the standard's
maintainers.

\paragraph{Scope.}
We stress what this paper is and is not. It is a \emph{feasibility and
identifiability audit}: we characterise, as a function of standard compliance,
when OD reconstruction is possible at all, and we validate the estimator's
convergence on synthetic data where the true cost is known. It is \emph{not} a
claim of validated pointwise OD reconstruction on the real feeds---no public feed
exposes both persistent identifiers and ground-truth trips---so the real data
serve only to instantiate the bound's constants $(q,\kappa,\beta,\gamma,\Lambda)$
and to place each network on the feasibility map. Because the novelty rests on
\emph{not} re-deriving the classical estimators, we first delimit what the
existing two literatures already settle.

% =========================================================================
\section{Background and related work}
\label{sec:related}

\paragraph{Counts give an underdetermined, prior-driven OD.}
Recovering an OD matrix from link or node counts is a 45-year-old, well-posed
but underdetermined problem. \citet{VanZuylenWillumsen1980} introduced the
maximum-entropy and minimum-information estimators; \citet{Cascetta1984}
incorporated survey priors through generalized least squares;
\citet{Spiess1987} gave the Poisson maximum-likelihood form; and
\citet{CascettaNguyen1988} unified the family. A practical point that frames
our own method is that the workhorse estimators are \emph{the same object}: the
iterative proportional fitting used to balance a matrix to target margins is
exactly the Sinkhorn iteration for entropic optimal transport
\citep{Sinkhorn1964,Cuturi2013,PeyreCuturi2019}, which is in turn the
maximum-entropy gravity model of \citet{Wilson1967}. Modern work sharpens the
conditioning of the inverse \citep{MaQian2023} but does not change the verdict:
from counts alone, pair-level OD is weakly identified.

\paragraph{Identifier tracking is exact but treats its biases as nuisances.}
When a feed exposes a persistent identifier, trips are reconstructed by tracking
position jumps between snapshots; \citet{Xu2022} give the reference algorithm
family, validated on one week of feeds from six vendors in a single city. This
and adjacent studies catalogue the biases---rebalancing contamination, GPS
jitter, and aliasing under coarse polling---but treat them as cleanup steps, and
are almost universally single-city, single-window, and silent on censored
demand \citep{RudloffLackner2014,Negahban2019,Goh2019} and on the standard's
identifier-rotation guidance.

\paragraph{The gap.}
No prior work models the impact of a \emph{standard specification}---concretely,
the identifier-rotation (persistence) rate $q$---on the mathematical solvability
of the downstream OD model. The closest cross-domain analogue, network
tomography \citep{Vardi1996}, draws its identifying power from a known routing
matrix that the station-count setting simply does not possess, so it collapses
to bare margin balancing. We close this gap by making compliance a parameter of
the estimator and bounding identifiability and collection cost as functions of
it. We now fix notation and the structural assumption that makes this possible.

% =========================================================================
\section{Problem formulation and data regimes}
\label{sec:formulation}

Let stations be indexed $i,j\in\{1,\dots,N\}$ and let $\Pstar\in\R_{\ge
0}^{N\times N}$ be the true OD coupling over a fixed period, with row margins
$r^\star_i=\sum_j \Pstar_{ij}$ (outflow) and column margins
$c^\star_j=\sum_i \Pstar_{ij}$ (inflow). Throughout we adopt the Gibbs
(gravity / entropic optimal-transport) structural assumption
\begin{equation}
\label{eq:gibbs}
\Pstar_{ij} = u^\star_i\, v^\star_j\, e^{-\cstar_{ij}/\varepsilon},
\end{equation}
with potentials $u^\star,v^\star$, ground cost $\cstar$, and temperature
$\varepsilon>0$. As recalled in Section~\ref{sec:related}, this is simultaneously
the maximum-entropy trip-distribution model and the fixed point of Sinkhorn
balancing, so \eqref{eq:gibbs} is not an additional modelling commitment beyond
the estimators the field already uses---only an explicit one.

Two observation channels constrain different parts of \eqref{eq:gibbs}.

\paragraph{Macro channel: censored margins.}
Station-status counts provide noisy, right-censored margins
$\hat r_i = \rho_i\, r^\star_i$ with availability fraction $\rho_i\in(0,1]$,
reflecting demand lost at empty or full stations \citep{RudloffLackner2014,
Negahban2019}. The macro channel thus pins (a degraded version of) the margins
but says nothing about the interior.

\paragraph{Micro channel: a selection-tilted sample.}
A sub-sample of identifier-tracked trips is drawn not from $\Pstar$ but from the
\emph{selection-tilted} law $\tilde P_{ij}\propto \Pstar_{ij}\,S_{ij}$, where the
selection probability $S_{ij}\in[0,1]$ aggregates identifier persistence $q$,
non-rebalancing $(1-\beta)$, and capture under polling interval $\Delta$
(anti-aliasing) $\kappa$. The micro channel thus carries information about the
interior---the cost---but through a biased lens. The remainder of the paper
asks what that lens lets us recover, and at what cost in collection time.

% =========================================================================
\section{Cost identifiability and the structure of observation bias}
\label{sec:theory}

\subsection{The transport cost is identifiable only modulo station effects}
\label{sec:identif}

The Gibbs form \eqref{eq:gibbs} is invariant under
$\cstar_{ij}\mapsto \cstar_{ij}+f_i+g_j$, since any additive separable term is
absorbed into the potentials. The cost is therefore identifiable only modulo the
separable subspace $\Ncal=\{(f_i+g_j)_{ij}\}$ ($\dim\Ncal=2N-1$); the
identifiable object is the interaction part $\Pperp \cstar$, equivalently the
second-order cross-differences $c_{ij}-c_{il}-c_{kj}+c_{kl}$.

\begin{theorem}[Cost identifiability]
\label{thm:identif}
Under \eqref{eq:gibbs}, $\Pperp\cstar$ is identifiable from a coupling, and no
separable component of $\cstar$ is. \Todo{State regularity conditions; the proof,
a matching/ground-metric argument in the manner of \citet{GalichonSalanie2022,
DupuyGalichon2014,CuturiAvis2014}, is given in SI Appendix~A.}
\end{theorem}

This is not a technicality but a design rule for the cost model. A low-rank
parametric cost $c_{ij}=\phi(x_i,x_j;\theta)$ may only use features whose
information survives projection onto $\Ncal^\perp$.

\begin{definition}[Admissible cost feature]
\label{def:admissible}
A candidate feature $\psi_{ij}=\phi^{(k)}(x_i,x_j)$ is \emph{admissible} iff it
survives two-way (origin/destination) centring,
$\tilde\psi_{ij}=\psi_{ij}-\bar\psi_{i\cdot}-\bar\psi_{\cdot j}+\bar\psi_{\cdot\cdot}
\neq 0$. Monadic features ($\psi_{ij}=a_i$ or $b_j$) are inadmissible: they are
absorbed into the potentials and belong in the margins, not the cost.
\end{definition}

Admissible, mutually non-collinear dyadic features---routed cycling distance,
\emph{directed} elevation gain, circuity, protected-lane fraction, and
river/rail barrier crossings---are what give the projected information matrix a
well-conditioned smallest eigenvalue $\lambda_{\min}$, which
Section~\ref{sec:horizon} shows is what keeps the collection horizon finite.

\subsection{Observation bias equals the non-separable log-selection}
\label{sec:bias}

Because the micro-sample is drawn from $\tilde P\propto \Pstar\odot S$ rather
than $\Pstar$, fitting \eqref{eq:gibbs} to it recovers the cost of $\tilde P$,
namely $\tilde c = \cstar - \varepsilon\log S$.

\begin{proposition}[Bias = non-separable log-selection]
\label{prop:bias}
The estimated identifiable cost satisfies
\begin{equation}
\label{eq:bias}
\Pperp\hat c \;\xrightarrow{\,n\to\infty\,}\; \Pperp\cstar
\;-\;\varepsilon\,\Pperp\log S .
\end{equation}
In particular, if the selection is separable, $S_{ij}=a_i b_j$, then
$\log S\in\Ncal$ and the bias vanishes. \Todo{Proof in SI Appendix~A.}
\end{proposition}

\subsection{Station emptiness cancels; 60-second aliasing does not}
\label{sec:aliasing}

Proposition~\ref{prop:bias} converts a question about data quality into a
question about the \emph{separability} of the selection mechanism, and the two
dominant defects fall on opposite sides of that line.

\begin{corollary}[Which censoring hurts]
\label{cor:aliasing}
Station-emptiness censoring acts at the origin/destination level
($S_{ij}=a_i b_j$) and therefore \emph{cancels} in the identifiable cost.
Aliasing under polling interval $\Delta$ has a capture probability that depends
on the trip duration $t_{ij}$, hence on $\cstar_{ij}$ itself; it is
\emph{non-separable} and correlated with the estimand, so it does not cancel and
systematically attenuates the cost--distance relationship. Residual operator
rebalancing, concentrated on specific corridors, is likewise non-separable.
\end{corollary}

The practical reading is uncomfortable: the defect one would expect to be fatal
(empty stations) is harmless to the cost, while the seemingly benign engineering
choice of a polling cadence is the true adversary, because it attacks exactly the
signal being estimated. Since the aliasing bias can only be beaten by sampling
more often and for longer, the natural next question is precisely \emph{how long}
one must collect---which leads us to formulate the sampling-horizon law.

% =========================================================================
\section{The sampling-horizon law}
\label{sec:horizon}

Let $\Lambda$ denote the trip-generation rate (trips per day). After $T$ days of
polling at interval $\Delta$, the number of clean, attributable observations is
\begin{equation}
\label{eq:neff}
\neff = \Lambda\,T\,\kappa(\Delta)\,q\,(1-\beta)\,(1-\gamma),
\end{equation}
where $\gamma$ is the GPS-jitter / collision loss. Propagating the cost-learning
error through the entropic-OT forward map---whose sensitivity is $O(1/\varepsilon)$
while its entropic bias is $O(\varepsilon)$, so that the optimal balance yields an
$n^{-1/4}$ rate on the OD error---turns a target accuracy into a required
effective sample $\neff\gtrsim \dim\theta/(\lambda_{\min}\,\delta_\OD^{4})$, and
hence a required collection time.

\begin{theorem}[Sampling-horizon law]
\label{thm:horizon}
To reconstruct the OD coupling to relative accuracy $\delta_\OD$ under the
estimator of Sections~\ref{sec:identif}--\ref{sec:aliasing},
\begin{equation}
\label{eq:Tstar}
\Tstar \;=\;
\frac{\dim\theta}
     {\lambda_{\min}\,\delta_\OD^{4}\;\Lambda\,\kappa(\Delta)\,q\,(1-\beta)(1-\gamma)} ,
\end{equation}
so $\Tstar\propto \delta_\OD^{-4}\,q^{-1}\,\Lambda^{-1}\,\lambda_{\min}^{-1}$.
\Todo{Forward-stability lemma and proof in SI Appendix~B.}
\end{theorem}

The fourth-power dependence on accuracy is the signature of the entropic
estimator: tightening the OD target from $20\%$ to $10\%$ multiplies the required
horizon by $2^4=16$. The inverse dependence on $q$ is the policy-relevant term.

\begin{corollary}[Standard-compliance threshold]
\label{cor:threshold}
For any operational horizon $T_{\max}$ and accuracy $\delta_\OD$, OD is
reconstructible only on feeds with identifier persistence
$q \ge q_{\min}:=\dim\theta/(\lambda_{\min}\delta_\OD^4 \Lambda
\kappa(1-\beta)(1-\gamma)\,T_{\max})$. Below $q_{\min}$ the horizon exceeds
$T_{\max}$ and operational monitoring is infeasible.
\end{corollary}

\begin{remark}[Stationarity ceiling]
\label{rem:stationarity}
Since $\Tstar$ is a \emph{cumulative} sampling requirement, it is only meaningful
below the OD process's stationarity period $T_{\mathrm{stat}}$; diurnal and
weekly cycles put $T_{\mathrm{stat}}$ at the order of a week for a static cost.
The operative constraint is therefore $\Tstar \le \min(T_{\max},
T_{\mathrm{stat}})$: if $\Tstar > T_{\mathrm{stat}}$ the feed mutates before
enough data accrue, so the coupling is formally unidentifiable within any single
regime. Non-stationarity thus tightens, rather than relaxes, the infeasibility
verdict.
\end{remark}

\subsection{Pooling is a forced trade-off, not a free lunch}
\label{sec:pooling}

For a single small network, \eqref{eq:Tstar} can run to years (Section~\ref{sec:audit}),
and the only term the analyst controls is $\dim\theta$. Modelling the cost
hierarchically---a deterrence law shared across the $N$ networks with city-level
random effects---divides the per-city parameter budget by the number of poolable
cities, and is the only mechanism that brings $\Tstar$ from ``years'' back to
``weeks''. We are careful not to present this as a free lunch: a shared $\theta$
assumes that distance, gradient and infrastructure deterrence are common across
cities, which they are not, so pooling buys a finite horizon at the price of a
specification bias that no individual city exactly satisfies. The honest reading
is that, to recover OD on small privacy-compliant feeds at all, the analyst is
\emph{forced} to sacrifice local specificity---itself a further indication that
the bottleneck lies in the standard, not in the estimator. A fixed $\varepsilon$
imposed for cross-city comparability additionally caps the achievable accuracy
at $\delta_\OD\gtrsim\varepsilon$. With the law in hand, we turn to measuring its
constants on real feeds.

% =========================================================================
\section{Auditing the French micromobility feeds}
\label{sec:audit}

\subsection{Protocol: measuring the five governing constants}
\label{sec:protocol}

The five constants $(q,\kappa,\beta,\gamma,\Lambda)$ are feed-specific and
\emph{measured}, not assumed. A polling daemon collects
\texttt{free\_bike\_status}/\texttt{vehicle\_status} snapshots at
$\Delta=60$\,s across the French free-floating systems that expose a vehicle
feed, writing one record per vehicle per poll (identifier, position, status
flags). Two reproducible scripts realise the estimators:
\texttt{experiments/d01\_pilot.py} computes the horizon \eqref{eq:Tstar} from a
set of constants, and \texttt{experiments/d02\_persistence\_audit.py} estimates
$\hat q$ feed by feed. Persistence is read off the snapshot timeline by rank over
the observed poll grid---robust to missed collection polls, which a clock-based
rule would miscount as trips: a vehicle whose identifier vanishes for one or more
served polls and \emph{reappears} is a linked disappearance, one that vanishes
before the window's safe sub-grid and never returns is terminal, and
$\hat q = \text{linked}/(\text{linked}+\text{terminal})$. The aliasing curve
$\kappa(\Delta)$ is obtained by re-reconstructing at coarser intervals,
rebalancing $\beta$ from the operator status flags and departure bursts, and
$\Lambda$ from the reconstructed clean trips. \Todo{Tabulate the full per-feed
constant table and the cost-feature pipeline (routing, elevation, infrastructure
layers) here; software stack and version pins; master seed $\texttt{SEED}=42$.}

\subsection{The compliance--horizon landscape}
\label{sec:landscape}

\begin{figure}[t]
  \centering
  \figorplaceholder{fig_d01_horizon}{0.82\linewidth}{Money plot ($T^\star$ vs
  $q$, from \texttt{d01\_pilot.py}): collection horizon against
  identifier-persistence rate for several trip rates $\Lambda$, with the
  operational-monitoring frontier and the audited feeds positioned on it.}
  \caption{The compliance--horizon landscape. The required collection horizon
  $\Tstar$ diverges as $q\to 0$ and as $\delta_\OD$ tightens (power $-4$);
  feeds below the operational frontier are not monitorable in practice.
  \Todo{regenerate with the multi-day constants.}}
  \label{fig:horizon}
\end{figure}

\begin{figure}[t]
  \centering
  \figorplaceholder{fig_d02_persistence}{0.82\linewidth}{Per-feed
  identifier-persistence $\hat q$ (from \texttt{d02\_persistence\_audit.py}),
  sorted, with the rotating/persistent thresholds marked.}
  \caption{Identifier persistence is bimodal across French feeds: $\hat q$ is
  either $\approx 0$ (standard-compliant, untrackable) or $\approx 1$
  (non-compliant, fully trackable), with little mass in between.}
  \label{fig:persistence}
\end{figure}

A first-light audit of $27$ French free-floating networks, over a short pilot
window with collection ongoing, already exhibits the split the theory predicts.
Of the $17$ feeds on which $q$ is estimable, $11$ \emph{rotate} their identifiers
($\hat q\approx 0$: vehicles vanish on rental and never return under the same
identifier---e.g.\ the largest free-floating fleet shows $0$ linked against
$>2000$ terminal disappearances), and $6$ \emph{persist}
($\hat q\in[0.90,1.00]$); the remaining $10$ are undetermined over the pilot
window (low-turnover car-sharing and very small fleets). The distribution is
sharply bimodal (Figure~\ref{fig:persistence}): $\hat q$ clusters at the two
extremes, essentially nothing in between, which is exactly the signature of a
binary engineering choice---whether or not an operator implements the standard's
rotation guidance---rather than a continuum of behaviours.
\Todo{Replace the pilot-window counts with the multi-day audit once collection
matures; report the full constant table and place each network on
Figure~\ref{fig:horizon}.} Read through Corollary~\ref{cor:threshold}, the
rotating half of the panel is simply not OD-monitorable from GBFS at any
practical horizon---a conclusion that is about the standard, not the operators,
and that we now develop.

% =========================================================================
\section{Discussion: standards and policy implications}
\label{sec:discussion}

Corollary~\ref{cor:threshold} turns a privacy feature of the standard into a
quantitative limit on what the standard can be used for. GBFS issue~\#146 and the
v2.0 guidance recommend rotating the vehicle identifier after each trip precisely
to prevent the full-resolution OD reconstruction that a persistent identifier
enables; the same property, by Theorem~\ref{thm:horizon}, is what drives the
collection horizon to infinity. There is therefore an irreducible tension, which
we make explicit rather than paper over: a feed cannot be simultaneously
privacy-compliant (small $q$) and OD-monitorable at a practical horizon.

For the standard's maintainers (MobilityData) this yields a concrete
recommendation. If GBFS is to remain the interface through which agencies expect
to derive mobility statistics, it must either (i) specify a minimum, auditable
persistence regime---e.g.\ session-scoped but trip-stable identifiers---or
(ii) state plainly that OD is \emph{out of scope} for a privacy-compliant GBFS
and route transactional needs to the regulator-facing Mobility Data
Specification (MDS), whose two-way, trip-level access is the appropriate channel
when $q<q_{\min}$. \Todo{Develop the $q_{\min}$ figure for representative French
agencies and connect to the LOM/GDPR context.}

One might object that regulators do not need GBFS at all, since MDS already
carries trip-level data. The objection sharpens our point. MDS access is
restricted to the operator's contracting authority, whereas GBFS is the only
\emph{public} channel available to researchers, civil-society groups, and the
many small municipalities without an audit mandate. If a strictly
privacy-compliant GBFS (small $q$) becomes unusable for macroscopic monitoring,
the operator--public information asymmetry becomes total for precisely those
actors, and Corollary~\ref{cor:threshold} quantifies the price of that
asymmetry. Either way, the design choice should be made knowingly, with the
horizon law in view, rather than emerging by accident from an
identifier-rotation default.

% =========================================================================
\section{Conclusion and limitations}
\label{sec:conclusion}

We have shown that the feasibility of OD reconstruction from GBFS is governed by
the standard's identifier-persistence regime: the learnable cost is identifiable
only up to station effects, its observation bias is exactly the non-separable
part of the log-selection (so station emptiness is harmless but 60-second
aliasing is not), and the resulting estimator obeys a sampling-horizon law
$\Tstar\propto\delta_\OD^{-4}q^{-1}$ that defines a compliance threshold below
which monitoring is infeasible. An audit of French feeds confirms a sharp
compliant/non-compliant split. The paper thus models the \emph{impossibility} of
reconstruction under a privacy-compliant standard rather than proposing a
reconstruction method: even under the strongly structuring Gibbs assumption the
horizon is operationally infeasible for compliant feeds, and any relaxation of
that assumption only weakens identifiability further and loosens the bound.

The results are conditional on the Gibbs/low-rank-cost assumption, which---though
it coincides with the field's standard estimators---is the binding commitment;
the stability constants in Theorem~\ref{thm:horizon} can be loose at large $N$
\citep{MenaNilesWeed2019,Genevay2019}; and the estimator recovers the
entropy-optimal representative, not pair-level OD, whose interior remains weakly
identified. Validation against ground truth requires a feed that both exposes
persistent identifiers and publishes trips, a thin intersection we address in
future work alongside the multi-day audit and the hierarchical cost of
Section~\ref{sec:pooling}.

% =========================================================================
% Back matter -- iopjournal native commands.  Migrate into elsarticle
% \frontmatter (and a CRediT statement) at submission.
% =========================================================================

\ack{\Todo{Acknowledgements.}}

\funding{\Todo{Funding sources and grant numbers.}}

\roles{\Todo{CRediT role assignments per author.}}

\data{All data sources and per-experiment derived outputs are released under
the MIT licence at
\url{https://github.com/cycling-data-lab/gbfs-od-reconstruction}. A frozen
Zenodo snapshot of the v1.0 release tag is deposited at
\url{https://doi.org/10.5281/zenodo.XXXXXXX} \Todo{substitute the issued
Zenodo DOI after the first release}.}

\section*{Reproducibility}
Each table and figure is produced by a numbered script
\texttt{experiments/d<NN>\_<name>.py} that consumes a known input cache, writes a
named output JSON/CSV, and pins \texttt{SEED=42}. The persistence audit
(Section~\ref{sec:audit}) is \texttt{d02\_persistence\_audit.py}; the horizon
figure is \texttt{d01\_pilot.py}. \Todo{Add experiment-specific detail.}

\section*{Competing interests}
The authors declare no competing financial or non-financial interests.

% =========================================================================
\section*{References}
\bibliography{references/references}

\end{document}
